\documentclass[12pt,a4paper]{article}

\usepackage[margin=1.8cm]{geometry}
\usepackage{fontspec}
\setmainfont{Times New Roman}
\usepackage{polyglossia}
\setdefaultlanguage{ukrainian}

\usepackage{xcolor}
\usepackage{enumitem}
\usepackage{setspace}
\usepackage{parskip}
\usepackage{tabularx}
\usepackage{array}
\usepackage{booktabs}
\usepackage{microtype}

\setlength{\parindent}{0pt}
\setlength{\parskip}{4pt}
\setstretch{1.03}

\definecolor{main}{HTML}{1F3A5F}
\definecolor{accent}{HTML}{6F7D8C}
\definecolor{linegray}{HTML}{CCD3DA}
\definecolor{textgray}{HTML}{444444}
\definecolor{soft}{HTML}{F5F7F9}

\newcommand{\topruleline}{%
    {\color{linegray}\rule{\linewidth}{0.9pt}}
}

\newcommand{\midruleline}{%
    {\color{linegray}\rule{\linewidth}{0.6pt}}
}

\newcommand{\formsection}[1]{%
    \vspace{0.8em}
    {\large\bfseries\color{main} #1}\par
    \vspace{0.35em}
    \midruleline
    \vspace{0.45em}
}

\newcommand{\checkboxitem}[1]{\item[□] #1}

\newcommand{\answerline}{%
    \vspace{0.9em}
    \hrulefill
    \vspace{0.9em}
}

\newcommand{\twolines}{%
    \answerline
    \answerline
}

\begin{document}

\begin{center}
    {\footnotesize\color{accent}\textsc{Матеріали емпіричного дослідження}}\\[2pt]
    {\small\color{textgray}\textsc{Психолого-педагогічний аналіз дитячого медіаконтенту}}
\end{center}

\vspace{0.3em}
\topruleline

\vspace{0.9em}

\begin{center}
    {\Large\bfseries\color{main} АНКЕТА ДЛЯ ВИХОВАТЕЛІВ}\\[0.45em]
    {\normalsize\bfseries\color{textgray}
    щодо оцінювання дитячого медіаконтенту,\\
    створеного за участю технологій штучного інтелекту}
\end{center}

\vspace{0.9em}

\begin{tabularx}{\linewidth}{@{}>{\bfseries}p{4.5cm}X@{}}
Мета опитування: &
виявити ставлення вихователів до якості, психологічної безпечності та доцільності використання дитячого медіаконтенту, створеного за участю штучного інтелекту.\\[0.55em]

Цільова група: &
вихователі закладів дошкільної освіти.\\[0.55em]

Форма участі: &
анонімна, добровільна.\\[0.55em]

Інструкція щодо заповнення: &
уважно прочитайте кожне запитання, позначте потрібний варіант відповіді або коротко впишіть власну думку у відведених рядках.
\end{tabularx}

\vspace{0.8em}
\topruleline

\vspace{0.8em}

\textbf{Шановні колеги!}

Просимо Вас взяти участь в опитуванні, результати якого будуть використані виключно в узагальненому вигляді в межах курсового дослідження. Анкета не передбачає зазначення персональних даних і спрямована на вивчення професійної думки педагогів щодо сучасного дитячого медіаконтенту для дітей дошкільного віку.

\formsection{1. Загальні відомості}

\textbf{1. Ваш педагогічний стаж:}
\begin{itemize}[leftmargin=1.5cm, itemsep=2pt, topsep=4pt]
    \checkboxitem{до 3 років}
    \checkboxitem{3--10 років}
    \checkboxitem{11--20 років}
    \checkboxitem{понад 20 років}
\end{itemize}

\textbf{2. З якою віковою групою дітей Ви працюєте?}
\begin{itemize}[leftmargin=1.5cm, itemsep=2pt, topsep=4pt]
    \checkboxitem{ранній вік}
    \checkboxitem{молодша група}
    \checkboxitem{середня група}
    \checkboxitem{старша група}
    \checkboxitem{різновікова група}
\end{itemize}

\formsection{2. Досвід використання дитячого медіаконтенту}

\textbf{3. Чи використовуєте Ви у своїй роботі відео, аудіо або цифрові матеріали для дітей?}
\begin{itemize}[leftmargin=1.5cm, itemsep=2pt, topsep=4pt]
    \checkboxitem{так, регулярно}
    \checkboxitem{інколи}
    \checkboxitem{дуже рідко}
    \checkboxitem{не використовую}
\end{itemize}

\textbf{4. Чи звертаєте Ви увагу на походження дитячого контенту (створений людиною чи за участю ШІ)?}
\begin{itemize}[leftmargin=1.5cm, itemsep=2pt, topsep=4pt]
    \checkboxitem{так}
    \checkboxitem{частково}
    \checkboxitem{майже не звертаю уваги}
    \checkboxitem{ні}
\end{itemize}

\textbf{5. Які критерії якості дитячого медіаконтенту для Вас є найважливішими?}
{\small\color{textgray}(можна обрати кілька варіантів)}
\begin{itemize}[leftmargin=1.5cm, itemsep=2pt, topsep=4pt]
    \checkboxitem{зрозумілість сюжету}
    \checkboxitem{спокійний емоційний фон}
    \checkboxitem{відповідність віку дітей}
    \checkboxitem{мовна грамотність}
    \checkboxitem{виховний зміст}
    \checkboxitem{естетичність оформлення}
    \checkboxitem{відсутність агресивних або тривожних елементів}
    \checkboxitem{інше: \hrulefill}
\end{itemize}

\textbf{6. Як Ви вважаєте, чи може контент, створений за участю ШІ, бути корисним для дітей дошкільного віку?}
\begin{itemize}[leftmargin=1.5cm, itemsep=2pt, topsep=4pt]
    \checkboxitem{так}
    \checkboxitem{скоріше так}
    \checkboxitem{скоріше ні}
    \checkboxitem{ні}
    \checkboxitem{важко відповісти}
\end{itemize}

\textbf{7. Які можливі ризики такого контенту Ви вбачаєте?}
{\small\color{textgray}(можна обрати кілька варіантів)}
\begin{itemize}[leftmargin=1.5cm, itemsep=2pt, topsep=4pt]
    \checkboxitem{емоційне перевантаження}
    \checkboxitem{надто яскраві або хаотичні стимули}
    \checkboxitem{поверховий зміст}
    \checkboxitem{невідповідність віковим особливостям}
    \checkboxitem{формування страхів або тривожності}
    \checkboxitem{зниження якості мовлення}
    \checkboxitem{не бачу особливих ризиків}
    \checkboxitem{інше: \hrulefill}
\end{itemize}

\textbf{8. Наскільки важливо, на Вашу думку, перевіряти дитячий контент перед показом дітям?}
\begin{itemize}[leftmargin=1.5cm, itemsep=2pt, topsep=4pt]
    \checkboxitem{дуже важливо}
    \checkboxitem{скоріше важливо}
    \checkboxitem{не завжди потрібно}
    \checkboxitem{неважливо}
\end{itemize}

\textbf{9. Які платформи або джерела дитячого контенту Ви використовуєте найчастіше?}
{\small\color{textgray}(можна обрати кілька варіантів)}
\begin{itemize}[leftmargin=1.5cm, itemsep=2pt, topsep=4pt]
    \checkboxitem{YouTube}
    \checkboxitem{YouTube Kids}
    \checkboxitem{телевізійні канали}
    \checkboxitem{освітні платформи}
    \checkboxitem{месенджери / соціальні мережі}
    \checkboxitem{файли від колег / методичні матеріали}
    \checkboxitem{інше: \hrulefill}
\end{itemize}

\formsection{3. Професійні спостереження та оцінка}

\textbf{10. Чи помічали Ви, що діти по-різному реагують на різні типи медіаконтенту?}
\begin{itemize}[leftmargin=1.5cm, itemsep=2pt, topsep=4pt]
    \checkboxitem{так, часто}
    \checkboxitem{інколи}
    \checkboxitem{майже не помічала(-в)}
    \checkboxitem{ні}
\end{itemize}

\textbf{11. Які реакції дітей Ви спостерігали найчастіше після перегляду медіаконтенту?}
\twolines

\textbf{12. Які ознаки, на Вашу думку, свідчать про психологічну безпечність дитячого контенту?}
\twolines

\textbf{13. Чи потрібні, на Вашу думку, рекомендації для педагогів і батьків щодо добору дитячого ШІ-контенту?}
\begin{itemize}[leftmargin=1.5cm, itemsep=2pt, topsep=4pt]
    \checkboxitem{так}
    \checkboxitem{скоріше так}
    \checkboxitem{скоріше ні}
    \checkboxitem{ні}
\end{itemize}

\textbf{14. Ваші додаткові думки або зауваження:}
\twolines

\vfill

\begin{center}
    {\color{linegray}\rule{0.36\linewidth}{0.8pt}}\\[0.45em]
    \textbf{\color{main}Дякуємо за участь в опитуванні!}
\end{center}

\end{document}